\begin{abstract}
This thesis investigates whether large language model (LLM) agents exhibit economically interpretable negotiation behaviour and whether coherent market-level dynamics emerge when many such agents interact under controlled bilateral bargaining conditions.  A configurable multi-agent simulator serves as the experimental instrument: it implements buyer and seller agents that negotiate through alternating-offer price bargaining with private economic constraints, and logs every agent action at sufficient granularity for behavioural analysis.  Six experiments (plus one extension) covering concession dynamics, deadline pressure, market price discovery, shock response, supply--demand shifts, informational anchor environments, and matching mechanism design were conducted across 3,420 sessions using \texttt{Qwen/Qwen2.5-14B-Instruct} with free-language agents and the feasibility gate disabled, so that every acceptance is LLM-driven.

At the negotiation level, LLM agents produce concave (Boulware-style) concession patterns with a 5.65:1 buyer-to-seller concession ratio. These results indicate that AI sellers maintain high price rigidity, sticking firmly to their initial anchors compared to their buyer counterparts.  Deadline pressure accelerates settlement: 82.4\% of deals close in the final two rounds under a six-round horizon.  

At the market level, prices exhibit strong stickiness in response to structural shocks (demand elasticity~0.275; supply elasticity~0.035).  An informational anchor environment experiment demonstrates that the reference price signal in the agent prompt is a primary determinant of price levels, with mean transaction prices spanning approximately \$30 across anchor modes under identical fundamentals.  Surplus-maximising matching improves total welfare by 7.2\% relative to random matching.  Taken together, these results reveal a micro--macro capability gap: agents close deals effectively at the session level yet fail to produce reliable price discovery in the aggregate.

Several behavioural patterns---concave concession trajectories, deadline-driven settlement acceleration, and reference-price sensitivity---are directionally consistent with patterns documented in the human negotiation literature, suggesting that LLM agents can serve as credible proxies for studying bargaining dynamics in controlled settings. The results establish a configurable experimental platform connecting micro-level agent behaviour to macro-level market outcomes, and lay groundwork for future work on prompt design, multi-model comparison, and institutional mechanism optimisation.

\end{abstract}
